\documentclass[11pt,a4paper]{article}
\usepackage[margin=1in]{geometry}
\usepackage{amsmath,amssymb,amsthm}
\usepackage{hyperref}
\usepackage{enumitem}

\newtheorem{theorem}{Theorem}
\newtheorem{definition}{Definition}
\newtheorem{axiom}{Axiom}
\newtheorem{proposition}{Proposition}
\newtheorem{lemma}{Lemma}
\newtheorem{remark}{Remark}
\newtheorem{criterion}{Criterion}

\title{MATHICCS: A Meta-Logical Criterion for Physical Theories\\
\large and Its Application to General Relativity}
\author{Steven E. Elliott}
\date{February 2026}

\begin{document}

\maketitle

\begin{abstract}
We introduce MATHICCS (Mathematics, Axiomatics, and Theory-Housed Internal Consistency
Constraints for Science), a formal meta-logical criterion for physical theories. The
criterion requires that every axiom of a physical theory possess an \emph{internal
epistemic realizer}: a physical apparatus, constructible from the theory's own
ontology, capable of executing the measurement procedure the axiom demands.

We apply this criterion to the Einstein Equivalence Principle (EEP), the foundational
axiom of General Relativity (GR). The EEP is an $\varepsilon$-$\delta$ statement about
spatial distances. GR itself defines spatial distance operationally via light travel
time. The minimum apparatus capable of executing the EEP's $\varepsilon$-$\delta$
distance procedure within GR's own framework is therefore a \emph{photonic mirror
clock}: a light pulse traveling a known distance $d$ and returning, defining the local
meter and second via the round-trip time $\Delta t = 2d/c$.

This apparatus is not an external philosophical imposition. It is the measurement
device GR's own operational definition of distance demands. We then show that GR
predicts physical conditions under which this device cannot be constructed or
maintained, and that those conditions are generic and persistent. The EEP therefore
becomes physically contentless across a substantial portion of the spacetime GR
describes, undermining GR's status as a complete physical theory of reality.

We further show that GR fails a stronger criterion, the \emph{Koons Hypothesis}: a
viable fundamental theory of reality must be such that evolution under its own dynamics
implies the structural stability of all its axioms. GR violates this because its own
dynamics destroy the apparatus required to give its foundational axiom physical content.

No measure-theoretic machinery is required for the operational argument. However, we
additionally demonstrate that the probabilistic form of the argument is equally
compelling: under GR's own predicted dynamics, the probability that any given
observer-moment falls within the thermodynamically active epoch is zero, meaning GR
assigns measure-zero probability to the conditions required for its own derivation.
GR cannot account for why its own discovery was possible without appealing to
something outside its own framework.
\end{abstract}

\tableofcontents
\newpage

%----------------------------------------------------------------------
\section{Introduction}
%----------------------------------------------------------------------

A physical theory does two things. First, it provides a mathematical formalism:
equations, geometric structures, transformation rules. Second, it asserts that this
formalism describes physical reality. The first task is purely mathematical; the second
requires a bridge between the formalism and observable phenomena. That bridge is
built from measurement.

Every foundational axiom of a physical theory carries an implicit measurement demand.
When GR asserts the Einstein Equivalence Principle, it is not merely making a statement
about the metric tensor in an abstract Riemannian manifold. It is asserting something
about how physical experiments behave in the real universe. The EEP has physical
content only because it predicts outcomes of measurements that can, in principle, be
performed.

The MATHICCS criterion asks a single question about any such axiom: does the theory's
own ontology support the existence of an apparatus capable of performing the required
measurement?

If yes, the axiom is \emph{internally realizable}: the theory is self-consistent in
the sense that it does not require external apparatus to give its axioms physical
meaning.

If no, the axiom is \emph{physically contentless} within the theory's own predicted
universe: the theory asserts something about measurement outcomes while predicting a
universe in which the relevant measurements cannot be performed. Such a theory has
a structural deficiency we call \emph{self-undermining}.

We demonstrate that GR is self-undermining in precisely this sense. The argument
requires no appeal to probability, no choice of measure, and no philosophical
controversy about the interpretation of physical theories. It requires only that we
take GR's own operational definition of distance seriously, identify the minimum
apparatus that definition demands, and ask what GR predicts about that apparatus.

We additionally show that GR fails the stronger \emph{Koons Hypothesis}
(Section~\ref{sec:koons}), which demands that a viable fundamental theory's own
dynamics must preserve, not destroy, the structural conditions required for its axioms
to have physical content. GR's dynamics actively dissolve those conditions.

The paper is organized as follows. Section~\ref{sec:mathiccs} introduces the MATHICCS
criterion. Section~\ref{sec:koons} states the Koons Hypothesis and its relationship to
MATHICCS. Section~\ref{sec:eep} identifies the EEP's measurement demand and minimum
realizer. Section~\ref{sec:destruction} shows GR predicts the destruction of that
realizer. Section~\ref{sec:theorem} states and proves the self-undermining theorem.
Section~\ref{sec:objections} addresses objections. Section~\ref{sec:implications}
discusses implications.

%----------------------------------------------------------------------
\section{The MATHICCS Criterion}
\label{sec:mathiccs}
%----------------------------------------------------------------------

\subsection{Physical Axioms and Their Demands}

\begin{definition}[Measurement Demand]
Let $A$ be a physical axiom of theory $\mathcal{T}$. The \emph{measurement demand}
of $A$, written $\mathcal{D}(A)$, is the set of physical operations that must be
performable in order for $A$ to have physical content beyond a purely mathematical
statement.
\end{definition}

\begin{definition}[Internal Epistemic Realizer]
Let $\mathcal{T}$ be a physical theory with ontology $\mathcal{O}$ (the set of
physical entities the theory posits as real). An \emph{internal epistemic realizer}
for axiom $A$ is a physical system $R$ such that:
\begin{enumerate}[label=(\roman*)]
\item $R$ is constructible from elements of $\mathcal{O}$ alone,
\item $R$ is capable of executing every operation in $\mathcal{D}(A)$, and
\item $R$ remains functionally intact under $\mathcal{T}$'s predicted dynamics for a
      duration sufficient to complete the required operations.
\end{enumerate}
\end{definition}

\begin{definition}[MATHICCS Validity]
\label{def:mathiccs}
A physical theory $\mathcal{T}$ is \emph{MATHICCS-valid} if and only if, for every
axiom $A$ of $\mathcal{T}$, an internal epistemic realizer for $A$ exists within
$\mathcal{T}$'s predicted universe.

A theory that possesses an axiom with no internal epistemic realizer is
\emph{MATHICCS-invalid}: it asserts physical content it cannot internally support.
\end{definition}

\begin{remark}
MATHICCS validity is a necessary condition for a theory to constitute a complete
physical description of reality. A MATHICCS-invalid theory may still be an excellent
effective tool within a restricted domain; but it cannot claim to be a fundamental
physical ontology, because it contains physical assertions that its own universe cannot
cash out in terms of measurement.
\end{remark}

\subsection{The Criterion is Theory-Internal}

A critical feature of Definition~\ref{def:mathiccs} is that the realizer must be
constructible from the theory's \emph{own} ontology. This is not an externally imposed
philosophical standard. It follows from the minimum requirement of scientific
realism: a theory that asserts ``the universe is such that measurement $M$ would yield
result $R$'' must, if taken as a description of reality rather than a conditional
mathematical statement, allow that measurement $M$ can actually be performed by
something inside that universe.

A theory that can only be validated from outside its own ontology is not a description
of physical reality. It is a mathematical structure that happens to fit some external
observations. The distinction matters when a theory makes claims about regimes where
external validation is impossible by construction.

\subsection{The Probabilistic Form of the Argument}
\label{sec:probabilistic}

Beyond the binary operational question of whether a realizer \emph{exists}, there is
a second, equally important way to state the deficiency. We include it here because it
answers a common objection --- ``why must realizers persist globally?'' --- and
illuminates the self-referential character of GR's failure.

The question is not whether laboratories must persist forever. The question is whether
a theory can account, from within its own dynamics, for why its own derivation was
possible at all.

If GR correctly describes the universe over $[t_{\mathrm{BB}}, +\infty)$, then a
randomly selected observer-moment in that universe has probability zero of falling
within the thermodynamically active epoch $[t_{\mathrm{BB}}, t_*]$. Einstein, any
working laboratory, any instance of GR being derived or tested --- all of these are
confined to a measure-zero slice of GR's own predicted history. GR therefore assigns
probability zero, under its own dynamics, to the conditions required for its own
discovery.

This is not merely a statement about improbability. It is a statement about
\emph{internal explanatory coherence}. A theory that assigns probability zero to
the conditions of its own derivation cannot explain, from within its own framework,
why those conditions obtained. It must either invoke an anthropic selection principle
(which GR's formalism does not contain), or concede that GR cannot account for its own
existence as a discovered theory without appealing to something outside itself.

Importantly, this argument is not defeated by restricting attention to the present
epoch. The objection ``we only need realizers now, and they exist now'' tacitly
assumes that the \emph{when} of derivation is given from outside the theory. But if
GR is the complete theory of the universe, there is no outside. The question of
\emph{when} derivation occurs must be answerable within GR's framework --- and GR's
framework says the probability of derivation occurring at any given time is zero in the
limit. One cannot help oneself to ``it happened to be now'' without invoking
something GR does not provide.

We note that this probabilistic argument is complementary to, not a replacement for,
the operational argument. The operational argument establishes that realizers are
absent in a definite region of GR's predicted spacetime. The probabilistic argument
establishes that GR cannot internally explain why derivation happened at all. Both
target the same structural deficiency from different angles.

%----------------------------------------------------------------------
\section{The Koons Hypothesis}
\label{sec:koons}
%----------------------------------------------------------------------

The MATHICCS criterion identifies a binary condition: a realizer either exists or it
does not. The \emph{Koons Hypothesis} is a stronger dynamical requirement, demanding
not merely that realizers exist but that the theory's own evolution preserves them.

\begin{criterion}[Koons Hypothesis]
\label{crit:koons}
Any viable or correct fundamental theory of reality must satisfy:
\[
\text{Evolution under its own dynamics} \;\Rightarrow\; \text{Structural stability of all its axioms.}
\]
That is: the dynamical laws of the theory must be such that the physical conditions
required to give the theory's axioms empirical content are preserved --- not destroyed
--- under the theory's own predicted evolution.
\end{criterion}

The Koons Hypothesis is stronger than MATHICCS validity in two respects. First, it is
a dynamical condition, not merely an existential one. MATHICCS asks whether a realizer
exists; the Koons Hypothesis asks whether the theory's dynamics \emph{maintain} that
existence. A theory could satisfy MATHICCS at early times while violating the Koons
Hypothesis by predicting the subsequent destruction of its own realizers. Second, the
Koons Hypothesis applies to the full temporal domain of the theory, not merely to some
initial epoch.

\begin{remark}
The Koons Hypothesis can be understood as a \emph{fixed-point condition} on the
theory's dynamics with respect to its own epistemic infrastructure. A theory satisfying
the Koons Hypothesis has dynamics that form a kind of closed loop: the laws govern
matter; matter instantiates the apparatus that gives the laws physical meaning; the
laws preserve that apparatus. A theory violating the Koons Hypothesis has laws that
destroy their own empirical grounding.
\end{remark}

\begin{remark}
The Koons Hypothesis does not require that realizers persist in every solution of the
theory. It requires that in the solutions the theory presents as physically realistic
--- in particular, solutions with the observed initial conditions --- the dynamics do
not generically destroy realizer existence. GR with $\Lambda > 0$ and observed initial
conditions is precisely the class of solutions under consideration.
\end{remark}

\subsection{GR and the Koons Hypothesis}

GR fails the Koons Hypothesis. Its dynamics do not merely fail to preserve its
axiom-realizers; they actively and necessarily destroy them. The mechanism is explicit:
de~Sitter expansion isolates all matter within a finite causal horizon; thermodynamic
processes dissipate all free energy on finite timescales; no mechanism within GR's
ontology regenerates the structural conditions required for a Photonic Mirror Clock
(Section~\ref{sec:eep}) after $t > t_*$.

This failure is not incidental. It is a consequence of GR's two most
observationally-confirmed features: the positive cosmological constant $\Lambda > 0$
(driving de~Sitter expansion) and the Second Law of Thermodynamics (driving irreversible
dissipation). GR's own best-confirmed predictions jointly entail the destruction of
its own realizers. The Koons Hypothesis failure is therefore not a contingent feature
of particular solutions; it is a structural property of GR as a cosmological theory.

Compare this to what a Koons-satisfying theory would require: its dynamics would need
to maintain, perpetually, the thermodynamic conditions necessary for its foundational
measurements. This is a genuine constraint on successor theories and points toward
theories with stationary, cyclic, or otherwise realizer-preserving dynamics.

%----------------------------------------------------------------------
\section{GR's Foundational Axiom and Its Measurement Demand}
\label{sec:eep}
%----------------------------------------------------------------------

\subsection{The Einstein Equivalence Principle}

\begin{axiom}[Einstein Equivalence Principle (EEP)]
\label{ax:eep}
For every point $p$ in the spacetime manifold $M$ and for every $\varepsilon > 0$,
there exists $\delta > 0$ and a coordinate system (local inertial frame) in which the
metric satisfies
\[
|g_{\mu\nu}(x) - \eta_{\mu\nu}| < \varepsilon \quad \text{for all } x \in U_\delta(p),
\]
where $\eta_{\mu\nu}$ is the Minkowski metric. In this frame, the laws of
special-relativistic physics hold to accuracy $\varepsilon$.
\end{axiom}

The EEP is the physical foundation of GR. It is what makes GR a theory about
the real world rather than a theory about abstract Riemannian manifolds. Without the
EEP, GR's field equations describe the evolution of a geometric structure with no
physical interpretation. With the EEP, the geometry is tied to the outcomes of
local measurements: rulers measure proper lengths, clocks measure proper times, and
local dynamics obey SR.

\subsection{The Measurement Demand of the EEP}

The EEP is an $\varepsilon$-$\delta$ statement. It asserts that spatial and temporal
distances can be measured locally to arbitrary precision $\varepsilon$ by shrinking the
neighborhood to size $\delta(\varepsilon)$. Its measurement demand $\mathcal{D}(\mathrm{EEP})$
therefore includes:

\begin{enumerate}[label=(\roman*)]
\item The ability to measure a spatial distance $d$ to precision $\varepsilon d$,
\item The ability to measure a time interval $\Delta t$ to precision $\varepsilon \Delta t$,
\item The ability to compare measurement outcomes at different points in $U_\delta(p)$,
\item The ability to detect deviations from SR physics at level $\varepsilon$.
\end{enumerate}

Operations (i) through (iv) must be executable by a physical system located within
the spacetime region under test. This is not optional: the EEP asserts something
about measurements performed at $p$, so the apparatus must be at $p$.

\subsection{GR's Own Definition of Distance Determines the Realizer}

Here is the key step, and it requires no external philosophical input.

GR defines spatial distance operationally. In GR, the proper distance between two
events is:
\[
ds^2 = g_{\mu\nu}\,dx^\mu\,dx^\nu.
\]
In a local inertial frame, this reduces to the SR line element. The operational
meaning of $ds$ in GR --- how one \emph{measures} a distance in the real world using
GR's framework --- is given by light travel time. Specifically, the proper length of
a spatial interval is defined as $c/2$ times the round-trip travel time of a light
signal across that interval. This is not a convention that can be replaced: in GR,
there is no distance-measuring procedure independent of light propagation. Rigid rods
are not primitive; they are defined via the metric, which is operationalized through
light.

Therefore, the minimum physical apparatus capable of executing the EEP's
$\varepsilon$-$\delta$ distance-measurement procedure is:

\begin{definition}[Photonic Mirror Clock (PMC)]
\label{def:pmc}
A \emph{Photonic Mirror Clock} is a physical system consisting of:
\begin{enumerate}[label=(\roman*)]
\item A source capable of emitting a directed light pulse,
\item A mirror at known separation $d$ from the source, capable of reflecting
      the pulse back to the source,
\item A detector at the source location capable of recording the round-trip
      arrival time $\Delta t = 2d/c$,
\item A record-keeping subsystem capable of storing and comparing multiple
      round-trip measurements.
\end{enumerate}
This device defines the local meter (via $d$) and the local second (via $\Delta t$),
and is capable of detecting deviations from the predicted $\Delta t = 2d/c$ at
precision $\varepsilon$.
\end{definition}

\begin{lemma}[PMC is GR's Own Minimum Realizer]
\label{lem:pmc}
The Photonic Mirror Clock is the minimum internal epistemic realizer for the EEP
within GR's own ontology. No simpler apparatus constructible from GR's ontology
can execute the EEP's $\varepsilon$-$\delta$ distance-measurement demand.
\end{lemma}

\begin{proof}
Any apparatus capable of measuring distance to precision $\varepsilon d$ in GR's
framework must operationalize $ds = g_{\mu\nu}dx^\mu dx^\nu$. In GR's operational
framework, this requires light propagation (as established above). The PMC is the
minimum system that accomplishes this: source, mirror, detector, and record-keeper.
Any attempt to reduce below this (e.g., remove the mirror) destroys the round-trip
measurement and with it the ability to define $d$ operationally. Any attempt to
remove the record-keeper destroys the ability to compare measurements across
$U_\delta(p)$ as required by the EEP. The PMC is therefore minimal.

We note that geodesic deviation --- the separation of neighboring free-fall
trajectories --- might appear to offer a non-photonic alternative. However, monitoring
the separation of two test masses requires measuring the distance between them, which
reduces back to light travel time in GR's operational framework. The geodesic deviation
apparatus is therefore a PMC under a different description, not an independent
alternative.
\end{proof}

\begin{remark}
We emphasize: the PMC is not an arbitrary experimental convenience. It is what GR
demands. A referee or critic who objects to the PMC as the realizer is implicitly
objecting to GR's own operational definition of distance. If one accepts GR's
framework, one accepts that distance is measured via light travel time, and therefore
one accepts that the PMC is the minimum realizer.
\end{remark}

%----------------------------------------------------------------------
\section{GR Predicts the Destruction of Its Own Realizer}
\label{sec:destruction}
%----------------------------------------------------------------------

We now examine what GR predicts about the existence and persistence of Photonic Mirror
Clocks in the universe GR describes.

\subsection{Requirements for PMC Existence}

A Photonic Mirror Clock requires:
\begin{enumerate}[label=(\alph*)]
\item \textbf{An energy source} to generate the light pulse. Photons carry energy
      $E = hf > 0$; generating them requires a thermodynamic process consuming free
      energy.
\item \textbf{A reflective mirror} at a stable, known separation $d$. This requires
      a bound material system --- atoms organized into a reflective surface --- held
      at fixed separation against perturbations.
\item \textbf{A detector} capable of registering photon arrival. This requires a
      material system with an accessible energy transition (absorbing or registering
      the photon).
\item \textbf{A record-keeper} capable of storing and comparing arrival times to
      precision $\varepsilon$. This requires a thermodynamically active system with
      sufficient complexity to encode and process information.
\end{enumerate}

All four components require \emph{thermodynamically active matter}: matter with
available free energy and organized structure at temperature above absolute zero
(and below a threshold that would thermally destroy the organization). This is a
single unified requirement:

\begin{definition}[Thermodynamic Viability]
A PMC is \emph{thermodynamically viable} at spacetime region $(x,t)$ if and only if
the matter at $(x,t)$ has sufficient free energy and structural organization to
construct and operate all four PMC components.
\end{definition}

\subsection{GR's Predictions About Thermodynamic Viability}

GR, together with the observed value of the cosmological constant $\Lambda > 0$ and
standard thermodynamics, predicts the following about thermodynamic viability in our
universe.

\subsubsection{Stellar Epoch}

Thermodynamically viable matter requires nuclear or chemical free energy. Stars provide
the dominant source. GR-consistent stellar physics predicts:

\begin{itemize}
\item Solar-mass stars: thermodynamic lifetime $\sim 10^{10}$ yr.
\item Minimum-mass hydrogen-burning stars: thermodynamic lifetime $\sim 10^{13}$ yr.
\end{itemize}

After stellar burnout, stellar remnants (white dwarfs, neutron stars, black holes)
cool monotonically and asymptote to zero temperature. They do not regenerate free
energy. The thermodynamic activity of the universe's matter therefore declines
monotonically after the stellar epoch ends.

\subsubsection{De Sitter Expansion and Isolation}

With $\Lambda > 0$, GR predicts exponential expansion of the scale factor:
$a(t) \propto e^{Ht}$ where $H = \sqrt{\Lambda/3}$. This produces a cosmological
event horizon of radius $r_H = c/H \approx 16$ Gly. As $t \to \infty$:

\begin{itemize}
\item All matter outside gravitationally bound local structures recedes beyond $r_H$
      and becomes causally inaccessible.
\item Within each isolated gravitational island, the matter supply is finite and
      fixed.
\item No thermodynamic replenishment from outside the horizon is possible.
\end{itemize}

\subsubsection{Baryonic Decay}

Grand Unified Theories consistent with GR predict proton decay on timescales
$\tau_p \sim 10^{32}$--$10^{40}$ yr~\cite{georgi1974}. Even absent proton decay,
quantum tunneling processes convert all matter to black holes and radiation on
timescales $\sim 10^{10^{76}}$ yr~\cite{adams1997}. Hawking radiation then
evaporates all black holes on timescales at most $\sim 10^{100}$ yr for stellar-mass
black holes. After these processes complete, no baryonic matter organized into
reflective surfaces, detectors, or record-keepers remains.

\subsubsection{Summary of GR's Prediction}

Let $t_*$ denote the latest time at which a thermodynamically viable PMC can exist
anywhere in the GR universe. By the above, $t_*$ is finite. An upper bound (extremely
conservative) is $t_* \lesssim 10^{10^{76}}$ yr. GR's predicted temporal domain is
$[t_{\mathrm{BB}}, +\infty)$.

Therefore:

\begin{proposition}[PMC Non-Existence]
\label{prop:pmc}
Under GR with $\Lambda > 0$ and standard thermodynamics, no Photonic Mirror Clock
exists in the GR universe for $t > t_*$, where $t_*$ is finite.
\end{proposition}

\noindent The universe GR describes spends the overwhelming majority of its predicted
existence in a state where no PMC can exist.

%----------------------------------------------------------------------
\section{The Self-Undermining Theorem}
\label{sec:theorem}
%----------------------------------------------------------------------

\begin{theorem}[GR Self-Undermining]
\label{thm:main}
General Relativity, taken as a complete physical theory of the universe with
$\Lambda > 0$, is MATHICCS-invalid and violates the Koons Hypothesis. Specifically:

\begin{enumerate}[label=(\roman*)]
\item \textup{(MATHICCS failure)} GR's foundational axiom (the EEP) has no internal
      epistemic realizer within GR's predicted universe for $t > t_*$, where $t_*$ is
      the finite epoch of thermodynamic viability. Since GR makes physical assertions
      about spacetime structure for all $t \in [t_{\mathrm{BB}}, +\infty)$, it makes
      physically contentless assertions about the entire post-thermodynamic domain.

\item \textup{(Koons Hypothesis failure)} GR's own dynamics --- de~Sitter expansion
      and thermodynamic dissipation --- do not preserve the structural stability of the
      EEP's realizer. They destroy it. The theory's evolution actively dissolves its
      own epistemic infrastructure.

\item \textup{(Probabilistic self-reference failure)} Under GR's own dynamics, the
      probability that a randomly selected observer-moment in the GR universe falls
      within the thermodynamically active epoch $[t_{\mathrm{BB}}, t_*]$ is zero.
      GR therefore assigns measure-zero probability to the conditions required for its
      own derivation and cannot internally account for why those conditions obtained.
\end{enumerate}
\end{theorem}

\begin{proof}
\textit{Part (i).} By Lemma~\ref{lem:pmc}, the minimum internal epistemic realizer for the EEP within
GR's ontology is a Photonic Mirror Clock. By Proposition~\ref{prop:pmc}, no PMC exists for $t > t_*$.
By Definition~\ref{def:mathiccs}, the EEP has no internal epistemic realizer for $t > t_*$.
Since GR asserts the EEP at all points of the spacetime manifold, including those with
$t > t_*$, GR asserts the physical content of the EEP in a regime where GR's own
ontology cannot support the realizer that would give the EEP physical content.
By Definition~\ref{def:mathiccs}, GR is MATHICCS-invalid.

\textit{Part (ii).} By Criterion~\ref{crit:koons}, GR must satisfy: evolution under
its own dynamics implies structural stability of all its axioms. GR's dynamics
(de~Sitter expansion with $\Lambda > 0$ plus thermodynamic dissipation) predict the
permanent destruction of all PMCs for $t > t_*$, where $t_*$ is finite and
independent of initial conditions within the observed cosmological class. The dynamics
do not preserve realizer existence; they terminate it. The Koons Hypothesis is
therefore violated.

\textit{Part (iii).} The temporal domain of GR is $[t_{\mathrm{BB}}, +\infty)$, which
has infinite Lebesgue measure. The thermodynamically active epoch
$[t_{\mathrm{BB}}, t_*]$ has finite Lebesgue measure $t_* - t_{\mathrm{BB}} < \infty$.
The fraction of observer-moments in a GR universe that fall within the active epoch is
$(t_* - t_{\mathrm{BB}}) / \infty = 0$. Under any measure on GR's temporal domain that
assigns positive weight to intervals of positive length in the post-thermodynamic
regime (which is required for $\mu$ to be a physically meaningful measure on an
unbounded domain), the realizer-supporting epoch has measure zero. GR therefore
assigns probability zero to the event ``derivation of GR is possible,'' and has no
internal mechanism to explain why derivation occurred.
\end{proof}

\begin{remark}
No step in part (i) invokes probability or measure theory. The argument is binary:
a PMC either exists in a given spacetime region or it does not. Part (iii) adds the
probabilistic form as an independent line of argument converging on the same
conclusion.
\end{remark}

\begin{remark}
The specific value of $t_*$ is irrelevant to parts (i) and (ii). Whether
$t_* = 10^{13}$ yr or $t_* = 10^{10^{76}}$ yr, the conclusion is the same: for
$t > t_*$, the PMC does not exist, the EEP is physically contentless, and the Koons
Hypothesis is violated. The specific value of $t_*$ matters only to the probabilistic
argument of part (iii), where any finite $t_*$ in an infinite domain yields measure
zero.
\end{remark}

%----------------------------------------------------------------------
\section{Objections and Responses}
\label{sec:objections}
%----------------------------------------------------------------------

\subsection{Objection 1: ``The EEP is a local axiom; it only needs to hold where observers exist.''}

\emph{Objection.} The EEP asserts that wherever an observer exists, local physics is
SR-like. If no observer exists, the EEP is vacuously satisfied, not violated. GR does
not assert that observers must exist everywhere.

\emph{Response.} A vacuously satisfied physical axiom is not a physical axiom. It is a
conditional mathematical statement: ``if an observer exists at $p$, then local physics
at $p$ is SR-like.'' This is weaker than the EEP as a physical postulate. GR's claim
is not merely conditional: GR asserts that spacetime has Lorentzian geometry, that
this geometry is dynamically determined by matter, and that local measurements reveal
SR physics. These are claims about the structure of the universe, not about the
conditional behavior of hypothetical observers.

More precisely: GR uses the EEP to derive the geodesic equation, the gravitational
redshift formula, and the equivalence of gravitational and inertial mass. These
derivations are used to make predictions about the present epoch. If the EEP is merely
a conditional statement, these derivations carry no physical force: they describe what
would happen if observers existed, but make no claim about whether observers do. This
would reduce GR from a theory of physical reality to a theory of observer-dependent
appearances. That is not what GR claims to be.

Furthermore, the Koons Hypothesis directly answers the ``local validity'' objection.
Even if one grants that the EEP need only hold locally where observers exist, the Koons
Hypothesis demands that the theory's dynamics maintain the conditions for observer
existence. GR's dynamics do not. A theory that is locally valid wherever observers
exist, but whose own dynamics eliminate all observers, is not locally valid
\emph{anywhere in its predicted future.} Local validity and dynamical stability of the
conditions for local validity are inseparable for a complete theory.

\subsection{Objection 2: ``Dyson showed that observers can persist indefinitely.''}

\emph{Objection.} Dyson (1979)~\cite{dyson1979} argued that life and computation can
persist indefinitely in an open universe by hibernation: slowing metabolism in
proportion to the decreasing ambient temperature, using finite total energy to survive
infinite time.

\emph{Response.} Dyson's argument applies to open universes without a cosmological
constant. With $\Lambda > 0$, Dyson's strategy fails for two independent reasons.

First, de~Sitter expansion imposes a \emph{finite} event horizon $r_H = c/H$. Any
observer is confined to a causally isolated region with finite total energy. Dyson's
hibernation strategy requires that available energy scales down proportionally to the
decreasing temperature; but with a finite energy reservoir and a nonzero Gibbons-Hawking
temperature~\cite{gibbons1977}
\[
T_{\mathrm{GH}} = \frac{\hbar H}{2\pi k_B} \approx 10^{-30}\,\mathrm{K},
\]
the observer must eventually expend more energy maintaining internal order against this
thermal bath than hibernation recovers. The finite energy budget is exhausted in
finite time.

Second, and more directly: even granting indefinitely persistent observers, a
Dyson-hibernating observer does not constitute a PMC. The PMC requires an active light
pulse, a reflective mirror, and a detector registering arrival time. A hibernating
observer by definition is not actively emitting and detecting photons. The hibernation
strategy preserves the \emph{existence} of a complex system but not its function as a
distance-measuring apparatus. The measurement demand of the EEP requires active
operation, not mere persistence.

The Koons Hypothesis sharpens this response further. A Dyson-hibernating observer is
not merely failing to operate a PMC; it is a system whose own dynamics are driving it
toward eventual dissolution against the Gibbons-Hawking bath. Even the hibernating
observer is on a finite trajectory under GR's dynamics. The Koons Hypothesis requires
the dynamics to maintain realizer structure indefinitely, not merely to delay its
destruction.

\subsection{Objection 3: ``This applies to any theory, not specifically GR.''}

\emph{Objection.} Any theory predicting eventual thermodynamic decay faces the same
argument. MATHICCS condemns all of physics, not just GR.

\emph{Response.} We accept this partially and consider it a feature, not a bug.
MATHICCS is a \emph{general} criterion, and the Koons Hypothesis is a general
requirement. Any theory predicting universal heat death and claiming physical validity
over its full temporal domain faces the same audit.

What MATHICCS demands is that a theory either (a) has a finite claimed domain over
which realizers persist throughout, or (b) explicitly restricts its physical validity
claims to the realizer-supporting epoch. Option (b) is available to GR: one can simply
say ``GR describes the thermodynamically active epoch, and makes no physical claims
about $t > t_*$.'' Most physicists implicitly make this restriction. MATHICCS makes it
explicit and mandatory.

The Koons Hypothesis adds a stronger requirement: a correct fundamental theory must
have dynamics that actively preserve its realizers. This points toward a genuine
structural requirement on successor theories: they must incorporate mechanisms
maintaining the thermodynamic conditions for their own foundational measurements.
Theories with cyclic, stationary, or otherwise realizer-preserving dynamics could in
principle satisfy the Koons Hypothesis; GR with $\Lambda > 0$ cannot.

\subsection{Objection 4: ``The PMC is too specific; other distance-measuring procedures exist.''}

\emph{Objection.} One could measure distances using gravitational wave detectors,
atomic clocks, or other devices not identical to a photonic mirror clock.

\emph{Response.} All such devices, examined carefully, reduce to the same operational
content: a signal (electromagnetic or gravitational) travels a known path and returns,
and the travel time defines the distance. Gravitational wave detectors measure the
differential light travel time in orthogonal arms --- they are PMCs. Atomic clocks
define frequency, not distance; to convert a frequency standard to a distance
measurement, one must still time a light signal.

The PMC in Definition~\ref{def:pmc} is not a specific experimental design; it is the
abstract minimum: a signal, a reflector, a detector, and a record-keeper. Any
distance-measuring device in GR's framework instantiates this abstract structure. The
argument applies to all of them equally.

\subsection{Objection 5: ``GR is not claimed to be a complete theory; everyone knows it needs quantum gravity.''}

\emph{Objection.} Physicists already regard GR as incomplete. MATHICCS is attacking a
straw man.

\emph{Response.} This is the most important objection and we address it directly.

It is true that GR is widely regarded as incomplete at quantum scales. But this
incompleteness is typically understood as a \emph{short-distance} or
\emph{high-energy} problem (singularities, quantization at the Planck scale).
MATHICCS identifies a different and independent incompleteness: a \emph{long-time},
entirely \emph{classical} problem. GR's physical axioms lose their realizers in the
asymptotic future, a regime that has nothing to do with quantum gravity and that
classical GR is taken to describe with full confidence.

These are orthogonal deficiencies. The standard incompleteness narrative runs:
singularities $\to$ quantum gravity needed $\to$ UV problem. MATHICCS identifies an
IR/asymptotic problem independent of quantization. GR could be successfully quantized
tomorrow and the MATHICCS failure would be entirely unaffected, because it concerns
the asymptotic classical future, not the quantum short-distance regime.

Furthermore, GR is routinely applied as a complete cosmological theory of the present
and future universe. Papers computing de~Sitter entropy, black hole information content
in the far future, or the thermal properties of empty de~Sitter space treat GR as
physically valid in those regimes. MATHICCS says: if you invoke GR's physical content
(the EEP) in those regimes, you must supply an internal realizer. If you cannot, you
are doing mathematics, not physics.

\subsection{Objection 6: ``Boltzmann brains could serve as realizers in the de~Sitter phase.''}

\emph{Objection.} Random thermal fluctuations in the de~Sitter phase could produce
Boltzmann brains that instantiate a PMC momentarily. These would constitute realizers
in the post-thermodynamic regime.

\emph{Response.} Boltzmann brains fail condition (iii) of the internal epistemic
realizer definition: they do not remain functionally intact under the theory's dynamics
for a duration sufficient to complete the required operations. A Boltzmann fluctuation
producing an apparent PMC exists only for the duration of the fluctuation; the same
thermal bath that produces it immediately begins to dissolve it. A Boltzmann PMC
cannot complete a measurement procedure, record results, and compare them, because it
does not persist as a functional system through the relevant sequence of operations.

More fundamentally: a Boltzmann-brain ``derivation'' of GR is not a legitimate
epistemic derivation. The Boltzmann brain does not arrive at GR through a causal
process of observation and inference; it merely instantiates a belief state that
matches GR. Such a system does not give the EEP empirical content; it merely contains
a representation of it. The epistemic realizer definition requires a system capable of
executing the measurement \emph{procedure}, not merely one that contains a record
asserting its results.

%----------------------------------------------------------------------
\section{What MATHICCS Does Not Claim}
\label{sec:limits}
%----------------------------------------------------------------------

For clarity, we enumerate what this paper does not assert.

\begin{enumerate}
\item We do not claim GR's predictions for $t \leq t_*$ are empirically wrong. GR's
remarkable agreement with observation within the thermodynamically active epoch is
unaffected by this argument.

\item We do not claim to identify the correct alternative theory. MATHICCS and the
Koons Hypothesis are filters, not replacements.

\item We do not claim the EEP is false. We claim it is physically contentless in
the post-thermodynamic regime, within GR's own ontology.

\item We do not invoke quantum mechanics or quantum gravity. The argument is classical:
GR plus standard thermodynamics.

\item We do not invoke probability or measure theory in the primary operational
argument. The primary claim is binary: a PMC either exists in a given spacetime region
or it does not. The probabilistic argument of Section~\ref{sec:probabilistic} and
Theorem~\ref{thm:main}(iii) is a complementary line of reasoning, not a replacement
for the operational argument.
\end{enumerate}

%----------------------------------------------------------------------
\section{Implications}
\label{sec:implications}
%----------------------------------------------------------------------

\subsection{The Minimum Requirement on Successor Theories}

Any theory that claims to supersede GR as a complete physical ontology must satisfy
both MATHICCS and the Koons Hypothesis. Concretely: it must predict that its own
minimum realizers persist throughout its claimed temporal domain, or explicitly
restrict its physical claims to the realizer-supported domain. Under the Koons
Hypothesis, the stronger requirement applies: the theory's dynamics must actively
maintain the conditions for its own epistemic grounding.

This is a structural constraint. It does not by itself determine what the successor
theory must look like. But it rules out any theory that, like GR, applies its
foundational axioms universally across a domain where those axioms have no internal
support, and any theory whose dynamics generically destroy their own realizers.

\subsection{The Self-Reference Character of the Problem}

There is a precise sense in which GR's deficiency is self-referential, visible from
three angles simultaneously.

Operationally: GR defines distance via light travel time; GR's dynamics eliminate all
light-travel-time measuring apparatus; GR therefore defines distance using a procedure
its own universe cannot perform.

Dynamically (Koons): GR's laws govern matter; matter instantiates GR's realizers; but
GR's laws destroy those realizers rather than preserving them.

Probabilistically: GR's temporal domain is infinite; the conditions for GR's discovery
occupy a finite slice; GR therefore assigns probability zero to its own discovery and
cannot explain, from within its own framework, why derivation occurred.

All three formulations target the same structural property: GR contains within its
formalism the prediction of its own epistemological dissolution.

\subsection{Connection to Known Problems}

The self-undermining property identified here is related to, but distinct from,
several known problems in foundations of physics:

\begin{itemize}
\item \textbf{The cosmological measure problem.} The measure problem asks how to
assign probabilities to observations in an eternally inflating universe. MATHICCS's
primary argument does not require probability: it asks only whether a realizer exists,
not how probable it is. The probabilistic argument of Section~\ref{sec:probabilistic}
intersects with the measure problem, but frames it as a disqualifying internal
inconsistency rather than an open technical puzzle.

\item \textbf{The Boltzmann brain problem.} MATHICCS does not rely on observer
counting. Boltzmann brains are excluded as realizers because they fail the functional
integrity condition of Definition~2, not because they are rare.

\item \textbf{The singularity problem.} GR predicts singularities where its own
equations break down at short distances. MATHICCS identifies an analogous breakdown at
long times: regions where GR's equations hold mathematically but their physical
interpretation is unsupported. The two problems are independent.

\item \textbf{The UV vs.\ IR split.} The standard incompleteness of GR is a UV (short
distance, high energy) problem. MATHICCS identifies an independent IR (long time, low
energy) problem. A successful UV completion of GR (quantum gravity) would not resolve
the MATHICCS failure.
\end{itemize}

%----------------------------------------------------------------------
\section{Conclusion}
\label{sec:conclusion}
%----------------------------------------------------------------------

We have established the following chain of reasoning, each step drawing on GR's own
framework:

\begin{enumerate}
\item The EEP is the foundational physical axiom of GR.
\item The EEP is an $\varepsilon$-$\delta$ statement requiring distance measurements
      at arbitrary precision.
\item GR defines distance operationally via light travel time.
\item Therefore the minimum internal realizer of the EEP within GR's own ontology is
      a Photonic Mirror Clock.
\item GR, together with standard thermodynamics and $\Lambda > 0$, predicts the
      permanent absence of PMCs for $t > t_*$, where $t_*$ is finite.
\item Therefore GR's foundational axiom has no internal realizer for $t > t_*$.
\item A physical theory whose foundational axiom has no internal realizer in a
      substantial portion of its claimed domain is not a complete physical description
      of reality. (MATHICCS failure.)
\item GR's dynamics do not preserve the structural stability of its axiom-realizers;
      they destroy them. (Koons Hypothesis failure.)
\item GR assigns probability zero, under its own dynamics, to the conditions required
      for its own derivation, and cannot internally account for why derivation
      occurred. (Probabilistic self-reference failure.)
\end{enumerate}

No step in this chain requires external philosophical commitments beyond the minimal
realist premise that physical theories must have internal support for their foundational
claims. The argument is self-contained within GR's own framework.

General Relativity, by its own predictions, destroys the apparatus required to give
its own foundational axiom physical meaning, violates the dynamical requirement that
a correct theory preserve its own epistemic infrastructure, and assigns measure-zero
probability to the conditions of its own discovery. This is what we mean by
MATHICCS-invalidity and Koons Hypothesis failure.

GR commits the same epistemic error as flat-earth theory: it postulates a globally valid geometric structure --- local flatness elevated to a universal axiom --- that cannot be tested from within the domain it claims to describe. Thus, Einstein created the 20th century version of flat earth.

\begin{thebibliography}{99}

\bibitem{bridgman1927}
P.~W.~Bridgman,
\textit{The Logic of Modern Physics},
Macmillan, New York (1927).

\bibitem{mach1883}
E.~Mach,
\textit{The Science of Mechanics},
Open Court (1883; English translation 1893).

\bibitem{gibbons1977}
G.~W.~Gibbons and S.~W.~Hawking,
``Cosmological event horizons, thermodynamics, and particle creation,''
\textit{Phys.\ Rev.\ D} \textbf{15}, 2738 (1977).

\bibitem{dyson1979}
F.~J.~Dyson,
``Time without end: Physics and biology in an open universe,''
\textit{Rev.\ Mod.\ Phys.} \textbf{51}, 447 (1979).

\bibitem{adams1997}
F.~C.~Adams and G.~Laughlin,
``A dying universe: The long-term fate and evolution of astrophysical objects,''
\textit{Rev.\ Mod.\ Phys.} \textbf{69}, 337 (1997).

\bibitem{georgi1974}
H.~Georgi and S.~L.~Glashow,
``Unity of all elementary-particle forces,''
\textit{Phys.\ Rev.\ Lett.} \textbf{32}, 438 (1974).

\bibitem{penrose2010}
R.~Penrose,
\textit{Cycles of Time},
Bodley Head, London (2010).

\bibitem{smolin2013}
L.~Smolin,
\textit{Time Reborn},
Houghton Mifflin Harcourt (2013).

\bibitem{norton2004}
J.~D.~Norton,
``The hole argument,''
\textit{Stanford Encyclopedia of Philosophy} (2004; revised 2019).

\bibitem{dyson2002}
L.~Dyson, M.~Kleban, and L.~Susskind,
``Disturbing implications of a cosmological constant,''
\textit{J.\ High Energy Phys.} \textbf{2002}(10), 011 (2002).

\end{thebibliography}

\end{document}